\chapter{Aug.~28 --- Convolution, Part 2}

\section{More Properties of Convolution}

\begin{exercise}
  Let $f, g \in L^1(\R)$, and denote
  $\widetilde{f}(x) = \overline{f(-x)}$.
  Showing the following:
  \begin{enumerate}[(a)]
    \item (Commutativity with translation)
      $f * (T_a g) = (T_a f) * g = T_a (f * g)$.
    \item (Behavior under involution)
      $(f * g)^\sim = \widetilde{f} * \widetilde{g}$.
  \end{enumerate}
\end{exercise}

\begin{remark}
  For $f, g : \R \to \C$, one can check that
  $f * g = 0$ outside
  \[
    \supp(f) + \supp(g)
    = \{x + y : x \in \supp(f), y \in \supp(g)\}.
  \]
  In particular, if $f, g \in C_c(\R)$, then
  $f * g \in C_c(\R)$ also.
\end{remark}

\section{Convolution and Smoothing}

\begin{theorem}\label{thm:conv-decay-at-infty}
  If $f \in L^1(\R)$ and $g \in C_0(\R)$, then
  $f * g \in C_0(\R)$.
\end{theorem}

\begin{proof}
  Since $C_0(\R) \subseteq L^\infty(\R)$, by
  Exercise \ref{ex:conv-cont-bounded} we have
  $f * g \in C_b(\R)$. To see the continuity
  more directly in this case, note that
  $g \in C_0(\R)$ implies $g$ is uniformly continuous.
  Thus for $x, h \in \R$,
  \begin{align*}
    |(f * g)(x) - (f * g)(x - h)|
    &= \left|
    \int f(y) g(x - y)\, dy - \int f(y) g(x - h - y)\, dy
    \right| \\
    &\le \int |f(y)| |g(x - y) - g(x - h - y)|\, dy \\
    &\le \bigg(
      \sup_{u \in \R} |g(u) - g(u - h)|
    \bigg) \int |f(y)|\, dy \\
    &= \|g - T_h g\|_\infty \|f\|_1
    \xrightarrow[h \to 0]{} 0
  \end{align*}
  by the uniform continuity of $g$.
  Alternatively, using the commutativity
  of translation with convolution,
  \[
    \|(f * g) - T_h(f * g)\|_\infty
    = \|f * g - f * T_h g\|_\infty
    = \|f * (g - T_h g)\|_\infty
    \le \|f\|_1 \|g - T_h g\|_\infty
  \]
  by Young's inequality. Either way, this
  establishes that
  $f * g \in C_b(\R)$.

  For the decay, assume first that $g \in C_c(\R)$,
  so $\supp(g) \subseteq [-N, N]$ for some $N > 0$.
  Then
  \[
    |(f * g)(x)|
    \le \int_{x - N}^{x + N} |f(y)| |g(x - y)|\, dy
    \le \|g\|_\infty \int_{x - N}^{x + N}|f(y)|\, dy
    \xrightarrow[|x| \to \infty]{} 0
  \]
  since $f \in L^1(\R)$.
  Thus $f * g \in C_0(\R)$ for any $g \in C_c(\R)$.

  For the general case, we extend by density.
  Given $g \in C_0(\R)$, by the density of
  $C_c(\R)$ in $C_0(\R)$, we can
  choose $g_n \in C_c(\R)$ such that $g_n \to g$
  in $C_0(\R)$.
  Then $f * g_n \in C_0(\R)$, and
  \[
    \|f * g - f * g_n\|_\infty
    \le \|f\|_1 \|g - g_n\|_\infty
    \xrightarrow[n \to \infty]{} 0
  \]
  Thus $f * g \to f * g_n$ in $L^\infty(\R)$.
  Since $C_0(\R)$ is a Banach space with
  $\|\cdot\|_\infty$, we
  get $f * g \in C_0(\R)$.
\end{proof}

\begin{remark}
  A somewhat shorter argument by density of
  Theorem \ref{thm:conv-decay-at-infty} goes as follows:
  Since $C_c(\R)$ is dense in both
  $L^1(\R)$ and $C_0(\R)$, we can choose
  $f_n, g_n \in C_c(\R)$ such that
  $f_n \to f$ in $L^1(\R)$ and
  $g_n \to g$ in $C_0(\R)$. Then
  as above, $f_n * g_n \in C_0(\R)$ for each
  $n$, and we can estimate
  \begin{align*}
    \|f * g - f_n * g_n\|_\infty
    &= \|f * g - f_n * g + f_n * g - f_n * g_n\|_\infty \\
    &\le \|f * g - f_n * g\|_\infty + \|f_n * g - f_n * g_n\|_\infty
  \end{align*}
  to get $f_n * g_n \to f * g$ in $L^\infty(\R)$.
  Then as above, we conclude that $f * g \in C_0(\R)$
  by completeness.
\end{remark}

\section{Convolution and Differentiation}

\begin{definition}
  Define the following spaces:
  \begin{align*}
    C_c^m(\R) &= \{f \in C_c(\R) : f, f', \dots, f^{(m)} \in C_c(\R)\}, \\
    C_0^m(\R) &= \{f \in C_0(\R) : f, f', \dots, f^{(m)} \in C_0(\R)\}.
  \end{align*}
\end{definition}

\begin{prop}
  Let $1 \le p < \infty$ and $m \ge 0$. If
  $f \in L^p(\R)$ and $g \in C^m_c(\R)$, then
  $f * g \in C_0^m(\R)$.
\end{prop}

\begin{proof}
  Assume first that $g \in C^1_c(\R)$. We can write
  \begin{align*}
    (f * g)'(x)
    &= \lim_{h \to 0} \frac{(f * g)(x) - (f * g)(x - h)}{h}
    = \lim_{h \to 0} \frac{1}{h} \int_\R f(y)(g(x - y) - g(x - h - y))\, dy \\
    &\overset{(*)}{=} \int_\R f(y) \lim_{h \to 0} \frac{g(x - y) - g(x - h - y)}{h}\, dy
    = \int_\R f(y) g'(x - y)\, dy
    = (f * g')(x),
  \end{align*}
  where the interchange of limit and integral in
  $(*)$ is justified by the dominated convergence theorem
  (find the dominating function as an exercise).
  Since $g \in C_c^1(\R)$, we have $g' \in C_c(\R)$,
  and so $(f * g)' = f * g' \in C_0(\R)$.
  The general case follows by induction (the above
  is essentially the inductive step).
\end{proof}

\begin{example}
  An example of a function
  $f \in C_c^\infty(\R)$ is
  \[
    f(x) = e^{-1 / x^2} \chi_{(0, \infty)}(x),
  \]
  which is smooth but not real analytic (taking
  $g(x) = f(x - a) f(b - x) \in C_c^\infty(\R)$ gives
  a bump function, which cannot be real analytic
  as vanishing on an open set would imply vanishing
  everywhere).
\end{example}

\section{Convolution on Other Domains}

\begin{definition}
  Given sequences $a = \{a_k\}_{k \in \Z}$
  and $b = \{b_k\}_{k \in \Z}$ of complex numbers,
  the \emph{convolution} is
  \[
    (a * b)_k = \sum_{j \in \Z} a_j b_{k - j},
  \]
  defined whenever the above series converges.
\end{definition}

\pagebreak

\begin{definition}
  Given $1$-periodic measurable functions
  $f, g : \R \to \C$, we can define their
  \emph{convolution}
  \[
    (f * g)(x)
    = \int_0^1 f(y) g(x - y)\, dy
  \]
  whenever the above integral is defined.
\end{definition}

\begin{remark}
  Given $1$-periodic
  $f, g : \R \to \C$, the
  convolution $f * g$ will also be $1$-periodic
  (if it is defined).
\end{remark}

\begin{remark}
  We can also view a $1$-periodic
  function $f : \R \to \C$ as a function
  $f : \mathbb{T} \to \C$, where $\mathbb{T} = [0, 1)$
  with addition defined modulo $1$.
  More generally, convolution can be defined in this way
  whenever one has a \emph{locally compact abelian group}
  $G$: Given a \emph{Haar measure} $\mu$ on $G$
  and $f, g : G \to \C$, we can define
  \[
    (f * g)(x) = \int_G f(y) g(x - y)\, d\mu(y).
  \]
  Every locally compact abelian group has a corresponding
  Fourier transform.
  One can also define a notion of convolution
  and Fourier transforms for nonabelian groups,
  though the theory is more difficult (there
  may be different
  left and right convolutions, and Fourier transforms may be matrix-valued in general).
\end{remark}
